\documentclass[8pt]{ctexbeamer}

\usepackage[chinese]{minimalistslides}

\setthemename{东南大学}
\setthemelogo{figs/logo.jpg}

\renewcommand{\presentationtitle}{大语言模型高效推理关键技术研究\\[1mm]\large 硕士学位论文开题答辩}
\renewcommand{\presentationauthorlabel}{汇报人}
\renewcommand{\presentationauthor}{付嘉乐}
\renewcommand{\presentationinstitutionlabel}{单位}
\renewcommand{\presentationinstitution}{东南大学}
\renewcommand{\presentationdate}{2026年4月}

\begin{document}

\makeboxedcover

\begin{frame}{目录}
    \tableofcontents
\end{frame}

\section{研究背景}

\subsection{背景与问题提出}
\begin{frame}{\insertsection}{\insertsubsection}
    \begin{columns}[T,onlytextwidth]
        \column{0.48\textwidth}
        \begin{block}{研究背景}
            \begin{itemize}
                \item 大语言模型已广泛应用于代码生成、数学推理与文本摘要等任务。
                \item 协同解码通过组合多个模型分布，通常可以提升生成质量与鲁棒性。
                \item 典型组合方式包括加权集成、Contrastive Decoding、Decoding-Time Realignment 等。
            \end{itemize}
        \end{block}

        \begin{block}{核心矛盾}
            质量提升依赖于多模型在每一步都参与预测，但这会显著增加推理延迟与算力开销。
        \end{block}

        \column{0.48\textwidth}
        \begin{figure}
            \centering
            \includegraphics[width=\linewidth]{ref/ref1/cos/figs/cod.pdf}
            \caption{协同解码在生成时融合多个模型分布。}
        \end{figure}
    \end{columns}
\end{frame}

\subsection{核心挑战与研究价值}
\begin{frame}{\insertsection}{\insertsubsection}
    \begin{columns}[T,onlytextwidth]
        \column{0.5\textwidth}
        \begin{block}{问题定义}
            多模型顺序调用意味着生成一个 token 往往需要多次前向计算，协同解码容易从“高质量”变成“高成本”。
        \end{block}

        \begin{block}{研究价值}
            本课题希望在\stronghighlight{不牺牲协同生成质量}的前提下，实现多模型推理加速，兼顾理论上的无损性与工程上的可部署性。
        \end{block}

        \column{0.46\textwidth}
        \begin{figure}
            \centering
            \includegraphics[width=\linewidth]{ref/ref1/cos/figs/efficiency_challenge.pdf}
            \caption{协同解码面临的效率瓶颈。}
        \end{figure}
    \end{columns}
\end{frame}

\section{研究内容}

\subsection{总体思路与核心科学问题}
\begin{frame}{\insertsection}{\insertsubsection}
    \begin{columns}[T,onlytextwidth]
        \column{0.52\textwidth}
        \begin{figure}
            \centering
            \includegraphics[width=\linewidth]{ref/ref1/cos/figs/sd_overall.pdf}
            \caption{推测解码的“先草拟后验证”范式。}
        \end{figure}

        \column{0.44\textwidth}
        \begin{block}{总体思路}
            以推测解码为切入口，将单模型的 proposal--verify 机制扩展到多模型协同解码场景。
        \end{block}

        \begin{block}{三个核心问题}
            \begin{enumerate}
                \item 协同分布如何参与验证并保持无损生成？
                \item 多个模型如何动态切换提议者与验证者角色？
                \item 该框架如何推广到 N 模型协同场景？
            \end{enumerate}
        \end{block}
    \end{columns}
\end{frame}

\subsection{研究内容一：协同分布下的推测机制}
\begin{frame}{\insertsection}{\insertsubsection}
    \begin{columns}[T,onlytextwidth]
        \column{0.52\textwidth}
        \begin{figure}
            \centering
            \includegraphics[width=\linewidth]{ref/ref1/cos/figs/ob1.pdf}
            \caption{将协同分布纳入验证过程的初步方案。}
        \end{figure}

        \column{0.44\textwidth}
        \begin{block}{研究目标}
            将传统推测解码中“目标模型分布验证”扩展为“协同/集成分布验证”，使输出分布严格对齐标准协同解码。
        \end{block}

        \begin{equation*}
            u \leq \min\left(1,\frac{\mathcal{C}(p(x_i), q(x_i))}{q(x_i)}\right)
        \end{equation*}

        \begin{block}{当前进展}
            已形成初步算法框架，后续将围绕拒绝采样准则、无损性证明和 acceptance rate 下界继续完善理论分析。
        \end{block}
    \end{columns}
\end{frame}

\subsection{研究内容二：动态验证与多模型推广}
\begin{frame}{\insertsection}{\insertsubsection}
    \begin{columns}[T,onlytextwidth]
        \column{0.5\textwidth}
        \begin{figure}
            \centering
            \includegraphics[width=\linewidth]{ref/ref1/cos/figs/ob22.pdf}
            \caption{交替提议框架示意图。}
        \end{figure}

        \column{0.46\textwidth}
        \begin{block}{动态角色分配}
            不固定“谁负责提议、谁负责验证”，而是依据上下文中的模型置信度与接受率，动态选择激活路径。
        \end{block}

        \begin{block}{多模型推广}
            从双模型扩展到 N 模型时，重点关注：
            \begin{itemize}
                \item 不同模型 proposal length 的协调；
                \item 加权集成场景下的加速比分析；
                \item 最坏情况下不慢于标准协同解码的性质。
            \end{itemize}
        \end{block}
    \end{columns}
\end{frame}

\subsection{研究方法与技术路线}
\begin{frame}{\insertsection}{\insertsubsection}
    \begin{block}{研究方法}
        \begin{enumerate}
            \item \textbf{理论推导：}基于概率论与拒绝采样，推导协同分布下的接受--拒绝准则与无损性条件。
            \item \textbf{算法实现：}在 HuggingFace / vLLM 框架上改造底层解码循环，构建多模型协同推测原型。
            \item \textbf{实验验证：}在代码生成、数学推理、文本摘要等任务上评估质量保持与速度提升。
        \end{enumerate}
    \end{block}

    \begin{block}{技术路线}
        \begin{enumerate}
            \item 复现标准协同解码与经典推测解码基线；
            \item 完成协同分布验证、动态角色切换与 N 模型推广设计；
            \item 引入 KV Cache 的 Lazy Synchronization / Forward Extend 机制，降低多模型切换带来的缓存维护开销。
        \end{enumerate}
    \end{block}
\end{frame}

\section{前期基础}

\subsection{阶段性结果与前期验证}
\begin{frame}{\insertsection}{\insertsubsection}
    \begin{table}
        \small
        \centering
        \renewcommand{\arraystretch}{1.25}
        \setlength{\tabcolsep}{3.2mm}
        \begin{tabular}{lccc}
            \toprule
            场景（代表性设置） & 并行实现 & 直接 SD & 阶段性结果 \\
            \midrule
            双模型 WE（Llama / Vicuna） & 0.69x--0.75x & 1.15x--1.27x & 1.41x--1.58x \\
            三模型 WE（Qwen-1.5B） & 0.54x--0.82x & 0.92x--0.98x & 1.27x--1.85x \\
            CD（Llama-3, $T=0$） & 0.40x--0.41x & 1.52x--2.04x & 1.61x--2.23x \\
            \bottomrule
        \end{tabular}
    \end{table}

    \begin{block}{阶段性观察}
        \begin{itemize}
            \item 任务覆盖 HumanEval、GSM8K、MMLU、CNNDM 等代表性基准。
            \item 初步结果表明：顺序协同虽能提升质量，但代价较高；简单并行实现未必更快。
            \item 协同推测框架在代表性设置下展现出 1.4x--2.2x 的潜在加速收益，说明技术路线具备可行性。
        \end{itemize}
    \end{block}
\end{frame}

\subsection{个人研究基础与相关成果}
\begin{frame}{\insertsection}{\insertsubsection}
    \begin{block}{与本课题最直接相关的前期工作}
        \textbf{ICML 2025}\\
        \textit{Fast Large Language Model Collaborative Decoding via Speculation}\\
        角色：第一作者；关联：直接支撑“协同解码 + 推测加速”的问题定义、方法设计与实验经验。
    \end{block}

    \begin{block}{相关研究积累}
        \begin{itemize}
            \item \textbf{ICLR 2026} \textit{Flatter Tokens are More Valuable for Speculative Draft Model Training}\\
            关联：为推测解码中的 draft model 训练与 acceptance behavior 提供理解。
            \item \textbf{ICLR 2026} \textit{$d^2$Cache: Accelerating Diffusion-Based LLMs via Dual Adaptive Caching}\\
            关联：为缓存机制设计与高效推理系统优化提供方法借鉴。
        \end{itemize}
    \end{block}

    \begin{myquote}
        更多成果可见：\texttt{jialefu.github.io/publications/}
    \end{myquote}
\end{frame}

\section{研究计划}

\subsection{研究计划、风险与预期成果}
\begin{frame}{\insertsection}{\insertsubsection}
    \begin{columns}[T,onlytextwidth]
        \column{0.5\textwidth}
        \begin{block}{阶段安排}
            \begin{table}
                \small
                \renewcommand{\arraystretch}{1.2}
                \setlength{\tabcolsep}{2mm}
                \begin{tabular}{p{0.24\textwidth}p{0.52\textwidth}}
                    \toprule
                    时间 & 主要任务 \\
                    \midrule
                    2026.04-06 & 完成文献梳理、基线复现与协同分布推导 \\
                    2026.07-09 & 完成动态验证策略与 N 模型框架设计 \\
                    2026.10-12 & 完成系统优化、消融实验与多任务评测 \\
                    2027.01-03 & 整理论文与学位论文，完善开源代码 \\
                    \bottomrule
                \end{tabular}
            \end{table}
        \end{block}

        \column{0.46\textwidth}
        \begin{block}{关键风险与应对}
            \textbf{风险：}多模型交替激活会带来 KV Cache 同步开销，可能抵消推测解码的加速收益。\\
            \textbf{应对：}采用 Lazy Synchronization 与 Forward Extend，仅在模型重新激活时批量补齐缺失缓存。
        \end{block}

        \begin{block}{预期成果}
            \begin{itemize}
                \item 形成可扩展的协同推测解码算法框架；
                \item 完成原型系统与核心代码整理；
                \item 在阶段性实验基础上推进论文与学位论文写作。
            \end{itemize}
        \end{block}
    \end{columns}
\end{frame}

\end{document}
